\chapter{Conclusion and Future Work}
\label{ch:conclusion}

\section*{Declaration of Contributions}
This chapter was jointly prepared by Parth Baranwal and Adesh Palkar,
with final review by all Group 4 members.

\section{Conclusion}

This report examined modern cyber threats and focused on software supply chain attacks
as a high-impact, trust-centric threat model. Through the XZ Utils case study
(\cref{ch:casestudy}) and related research synthesis (\cref{ch:background}), we observed that
a single compromised upstream dependency can propagate risk to many downstream systems.

The in-depth analysis (\cref{ch:solutions}) shows that existing solutions improve visibility
and control but remain individually insufficient. The central conclusion is that dependency
trust must be continuously verified, not implicitly assumed.

\section{Future Work}

The following directions are recommended for improving practical defenses:

\begin{enumerate}
\item standardized provenance and signed artifact verification across package ecosystems;
\item stronger CI/CD policy gates for dependency updates and transitive risk checks;
\item scalable hybrid analysis (static + dynamic + symbolic) with low false-positive rates;
\item faster ecosystem-wide notification and coordinated mitigation for exposed dependents;
\item broader adoption of security-by-default package publishing and maintainer hardening.
\end{enumerate}

These directions support a shift from reactive patching to proactive and continuous
software supply chain security.
